\section{The quotient-method system}
\label{sec:sec3}

Section~3 of \cite{jsww1976} contains a second, quite different construction,
following Matijasevi\v{c} and Robinson \cite{mr1975}: their Theorem~3.9 gives
$21$ conditions in which fourteen unknowns are then removed. The removal is
disposed of in one sentence (p.~461): ``The unknowns $M,A,B,C,D,E,F,G,H,I,K,L,R,S$
eliminate from (I)--(XXI) by substitution.''

Over $\N$ that sentence has content, by Proposition~\ref{prop:elim}. Applying
the structural criterion of Remark~\ref{rem:crit} licenses only six of the
fourteen; the remaining eight are each blocked by one subtraction. We prove all
eight, and formalize them.

\subsection{The eight bounds}

The fourteen unknowns split into six that pass the structural criterion
outright---$M$, $A$, $B$, $C$, $E$, $H$---and eight that do not: $D$, $F$, $G$,
$I$, $K$, $L$, $R$, $S$. It is those eight we must bound. Write
$U(X,Y)=(X+2)^3(X+4)(Y+1)^2+1$, so that conditions~(I) and~(II) of
Theorem~3.9 read $U(2k,n)=\sq$ and $U(2n,x)=\sq$.

Four of the eight---$D$, $F$, $G$, $I$---follow from an elementary chain. Since all variables are in $\N$,
$M=16nx(w+2)+1\ge1$, hence $A=M(x+1)\ge1$ and $A^2-1\ge0$---which is exactly the
subtraction blocking $D$ and $F$. Then $C=m+n+1\ge1$, $D\ge1$,
$E=2(i+1)DC^2\ge2$, and with $E\ge2$,
\[
  F=(A^2-1)E^2+1 \;\ge\; 4A^2-3 \;\ge\; A ,
\]
because $4A^2-A-3=(A-1)(4A+3)\ge0$; this is the subtraction blocking $G$, and
$G=A+F(F-A)\ge A\ge1$ unblocks $I$. Two more, for $L$ and $R$, need $Mx\ge1$ and
$Mnx\ge1$, hence $n\ge1$ and $x\ge1$; these follow from the observation that
$U(X,0)=t^4+2t^3+1$ with $t=X+2$ lies strictly between the consecutive squares
$(t^2+t-1)^2$ and $(t^2+t)^2$ and is therefore not a square, so (I) and~(II)
force exactly that.

The seventh, for $K$, requires $k\le n$. Condition~(I) reads
$16(k+1)^3(k+2)(n+1)^2+1=c^2$, which at first sight is not of the Pell form
$x^2-(A^2-1)y^2=1$ our development handles. It is, by the identity
\[
  16m^3(m+1) \;=\; \bigl((2m+1)^2-1\bigr)\cdot 4m^2 , \qquad m=k+1,
\]
so that $A=2k+3$ and $y=2(k+1)(n+1)$; the congruence then sends the divisibility
$2(k+1)\mid y$ to the index, giving $j\ge2(k+1)\ge4$, and growth from
$Y_3=4A^2-1$ yields $8k\le n+1$.

\subsection{The eighth bound was a transcription error of ours}

The last, for $S=(z+1)(k+1)-2$, requires $(z+1)(k+1)\ge2$, which fails at
exactly one point, $z=k=0$. It resisted proof because $z$ occurs \emph{only} in
that equation, so no other condition constrains it.

The resolution is that the point is not in the theorem. Theorem~3.9
of~\cite{jsww1976} begins ``For any \emph{positive} integer $k$'', and the
Wilson criterion it rests on (their Lemma~2.9) likewise requires $k\ge1$---as it
must, since at $k=0$ the criterion misfires: $k+1=1$ divides $k!+1=2$, so it
would report that $1$ is prime. We had transcribed the parameter as ranging over
$\N$, thereby asserting more than the theorem does. Under the exact
reparametrization of its domain, $k=k'+1$ with $k'\in\N$, the definition becomes
$S=k'z+k'+2z$, all of whose coefficients are nonnegative, and $S$ is eliminable
with no bound at all.

We record this because of what it turned out to mean: \emph{the point at which
our bound failed is exactly the point at which Wilson's criterion fails}. That
is not a coincidence. $S$ is the unknown that carries Wilson into the system---
condition~(XXI) says $k+1\mid S+2$, and the sufficiency proof uses $S\ge0$ to
conclude $\tfrac12<\beta$. The two were the same fact seen from two places.

\subsection{The inequality (XIV) carries no inherited hypothesis}

Condition~(XIV) of Theorem~3.9 asserts that a certain rational quantity
$\beta$ is within $\tfrac12$ of $S+1$, where
\[
  \beta \;=\; \frac{R}{\bigl(\tfrac{C}{KL}-(w+1)x\bigr)\bigl(1-\tfrac{R}{C}\bigr)^2 L}
        \;=\; \frac{\mathit{Nu}}{\mathit{De}}, \qquad
  \mathit{Nu}=RKC^2,\quad
  \mathit{De}=\bigl(C-(w+1)xKL\bigr)(C-R)^2 ,
\]
the second equality after clearing denominators (the factor $L$ cancels). Here
$K$, $L$, $R$, $C$ are among the fourteen unknowns and $w,x$ among the ten that
survive. Converting $|\beta-(S{+}1)|<\tfrac12$ into a polynomial equation
requires multiplying by $\mathit{De}^2$, which we had documented as inheriting
the hypothesis $\mathit{De}>0$ (formula~(15) of \cite{jsww1976}, p.~458). It
does not: the encoding forces it. The slack form
$4(\mathit{Nu}-(S{+}1)\mathit{De})^2+1+s_1=\mathit{De}^2$ with $s_1\ge0$ is
equivalent to $4(\mathit{Nu}-(S{+}1)\mathit{De})^2<\mathit{De}^2$, and

\begin{itemize}[nosep]
  \item $\mathit{De}=0$ gives $4\mathit{Nu}^2+1+s_1=0$, impossible over $\N$;
  \item $\mathit{De}<0$ gives
        $\mathit{Nu}-(S{+}1)\mathit{De}\ge1+|\mathit{De}|$---using
        $\mathit{Nu}\ge1$ and $S+1\ge1$---whence
        $4(\cdot)^2\ge4(1+|\mathit{De}|)^2>\mathit{De}^2$, contradicting the
        inequality.
\end{itemize}

Both hypotheses are among the eight bounds just proved: $\mathit{Nu}=RKC^2\ge1$
because $R,K,C\ge1$, and $S+1\ge1$ because $S\ge0$. Note that $S+1\ge1$ is
precisely what was lost at $k=0$: with $S+1=0$ the negative-denominator
solutions reappear. The two gaps closed together.

\subsection{The relation-combining theorem, and four published numbers}

Theorem~3 of \cite{mr1975} fuses $q$ square conditions, one divisibility and one
positivity into a single equation at the cost of one unknown. We implemented it
as an iterated norm (not by expanding $2^q$ factors) and validated the
implementation against four figures published in \cite{jsww1976}:

\begin{center}
\begin{tabular}{@{}llr@{}}
\toprule
\textbf{Quantity} & \textbf{Source} & \textbf{Value} \\
\midrule
$\deg M_6$, plain theorem & p.~461 & $148864$ \\
$\deg M_6$, with Matijasevi\v{c}'s refinements & p.~462 & $13376$ \\
$\deg M_5$, with condition (24) & p.~462 & $6848$ \\
generator degree & p.~462 & $13697$ \\
\bottomrule
\end{tabular}
\end{center}

Three of these are independent computations; the fourth, $13697=1+2\cdot6848$,
follows from the third by Proposition~\ref{prop:gen}. All are reproduced
exactly. The refinements are described telegraphically
on p.~462; the only reading not fixed by the text is the degree of the auxiliary
bounds $V_i$, for which we take $\lceil \deg A_i/2\rceil$---natural, since the
condition \emph{is} that $A_i$ be a square. Four exact agreements from a single
reading is our evidence that the reading is the right one.

Applying the refined theorem to subsets of the conditions gives a curve:
$(16,369)$, $(15,801)$, $(14,1777)$, $(13,3905)$, $(12,13697)$, whose last point
is the figure of \cite{jsww1976}. We are not aware of other pairs with fewer than $17$ variables and degree
below~$6000$, but we state this with two explicit reservations. First,
\cite{jsww1976} announce on p.~449 that ``the method of proof of Theorem~1
yields a polynomial in 16 variables'' without giving its degree; if that degree
is below $6000$, the claim fails, and we have not been able to determine it.
Second, these four points are \emph{announced by us}, not exhibited: we do not
write the polynomials down, they are not formalized, and their degrees are
\emph{upper} bounds computed by traversing the expression tree, since expanding
is infeasible. By our own criterion they are therefore in the same category as
the figures we criticise elsewhere in this paper, and we label them so.

They are also \emph{not} a record in the number of variables: the
$(10,>6000)$ of \cite{matiyasevich1981} has fewer variables than any of them.
We do not assert that it dominates the $(12,13697)$, because $>6000$ is a lower
bound and the comparison with $13697$ is therefore open.
